\section{Conclusions}
\label{sec:conclusion}
General-purpose engines apply fixed rewrite rules, LLM rewriters pay tokens for every new query or template, and engine synthesis bypasses the host engine. {\sys} teaches existing engines new rewrites for evolving workloads. The \emph{Offline Rule Generator} distills the output of any whole-query LLM rewriter into rules that are independent of constants, tables, and join arity. It admits a rule only after it is formally proven equivalent or passes result-equivalence tests. The \emph{Online Rewrite Engine} matches rules on plan fragments with no LLM call and emits SQL for the unmodified host system. Across nine workloads on PostgreSQL, MySQL, and DuckDB, {\sys} speeds up Stats-CEB by 38.7$\times$ on PostgreSQL versus 12.9$\times$ for ReSequel, transfers JOB rules unchanged to Stats-CEB for up to 28.3$\times$, and outperforms BespokeOLAP by 2.25$\times$ end-to-end.
%
We draw several major conclusions. First, structure transfers, not queries: on Stats-CEB, rules borrowed from JOB can outperform rules distilled from Stats-CEB itself. Second, the rule set saturates: 125 rules across seven workloads reduce to 67 distinct rules, so token cost follows new patterns, not new queries. Third, gains depend on what the host lacks: some rules are impactful on PostgreSQL and MySQL but neutral on DuckDB. Fourth, priority order is a strong default, and search helps only when rules compete at the same plan node. Fifth, verifying a rule once is more practical than verifying every rewritten query. Six, rewriting only the logical plan, and reusing existing execution engines, pays off: rewritten queries gain 5.5--7.8$\times$ from multi-threading, while a synthesized engine gains at most 1.69$\times$.